% Source: https://artofproblemsolving.com/community/c3h1552885_pythagorean_hypotnuses
% Stable identity: aops-topic-1552885-q01
\textbf{(a)} In a primitive triple, $a$ and $b$ are not both even (otherwise $2$ divides $a$, $b$ and $c$), and not both odd (then $c^2=a^2+b^2\equiv2\pmod4$, but squares are $0$ or $1$ modulo $4$). Say $a$ is odd and $b$ even; then $c$ is odd. Write
\[
b^2=c^2-a^2=(c-a)(c+a).
\]
Both factors are even; put $c+a=2u$ and $c-a=2v$ with positive integers $u,v$. Any common divisor of $u$ and $v$ divides $u+v=c$ and $u-v=a$, and then also $b$; so $\gcd(u,v)=1$. Now $\left(\dfrac {b}{2}\right)^2=uv$ with $u$, $v$ coprime, so each of $u$ and $v$ is a perfect square (their prime factorisations do not overlap, so every prime appears to an even power in each): $u=m^2$, $v=n^2$ with positive integers $m$, $n$. Hence
\[
c=u+v=m^2+n^2 .\qquad\blacksquare
\]
For example $5=1^2+2^2$, $13=2^2+3^2$, $17=1^2+4^2$, $25=3^2+4^2$.

\textbf{(b)} \textbf{No.} The triple $(9,12,15)$ (three times $(3,4,5)$) has hypotenuse $15$, and $15$ is not a sum of two squares: $15-0,15-1,15-4,15-9$, that is $15,14,11,6$, are not perfect squares. (Only non-primitive triples can fail, by part (a).)
